\documentclass[11pt]{article}
\usepackage[utf8]{inputenc}
\usepackage[T1]{fontenc}
\usepackage{lmodern}
\usepackage[margin=1in]{geometry}
\usepackage{amsmath,amssymb,amsthm,mathtools}
\usepackage{booktabs}
\usepackage{enumitem}
\usepackage[colorlinks=true,linkcolor=black,citecolor=black,urlcolor=black]{hyperref}
\usepackage{longtable}
\usepackage{array}
\usepackage{graphicx}
\usepackage{microtype}
\setlength{\parskip}{0.65em}
\setlength{\parindent}{0pt}

\title{Paper 5 --- REPRESENTATIONS (Quantum): Quantum Structure from Finite Enforceability}
\author{E. Brooke}
\date{}

\begin{document}
\maketitle
\begin{abstract}
This paper derives the mathematical structure of quantum mechanics---Hilbert space, the Born
 rule, tensor products, unitarity, and CPTP dynamics---from the admissibility framework established
 in Papers 1--4. Quantum structure is not postulated but emerges as the minimal bookkeeping
 required to represent coherent, jointly enforceable distinctions under finite capacity. The derivation
 proceeds without assuming probability, dynamics, or any quantum postulates.

 This February 2026 release incorporates the record-locking channel formalization (T9/L3-$\mu$), which
 provides an explicit algebraic model for measurement as irreversible record formation on M =
 M_sys $\otimes$ Z_R, yielding k! orthogonal history sectors. The connection to the derived gauge
 structure (T_gauge, T4F) is made explicit: quantum structure and matter content are jointly
 constrained by the same capacity architecture.
\end{abstract}

\section{Introduction}



\subsection{The Question}

Why does quantum mechanics have the mathematical structure it does? Standard presentations

take this structure as given: Hilbert spaces, complex amplitudes, the Born rule, tensor products,

unitary evolution. These are axioms or postulates, accepted because they work. But the formalism

permits vastly more structure than nature realizes. Why the specific subset that we observe?



\subsection{The Answer}

Quantum structure is not arbitrary. It is the minimal bookkeeping required to track admissible

distinctions under finite enforceability. This paper shows: (1) Linearity arises from coherence

(interference closure). (2) Inner products arise from admissibility-invariant overlap. (3) Complex

numbers arise from continuous reversible symmetry. (4) Tensor products arise from compositional

closure. (5) The Born rule is the unique admissibility-invariant probability measure. (6) Unitarity

preserves admissibility under reversible evolution. (7) CPTP maps preserve admissibility under

general evolution. No quantum postulates are assumed. The structure is forced by admissibility.



\subsection{What This Paper Assumes}



This paper assumes: Paper 1 (SPINE), Paper 2 (STRUCTURE), Paper 3 (LEDGERS), Paper 4

(CONSTRAINTS). It also assumes the bilinearity regime: the interaction functional I_$\Gamma$ is

approximately bilinear. This is a regime condition, not a universal truth.



\subsection{What This Paper Does NOT Assume}

Hilbert space (derived), complex numbers (derived), Born rule (derived), unitarity (derived),

measurement postulates (regime transitions, not postulates), any quantum axioms.



\section{Coherence and Linear Structure}



\subsection{Coherent Distinctions}

A set of distinctions is coherent if they can be reversibly refined, recombined, and undone without

violating admissibility. Coherence means distinctions can be made finer (refined), refined

distinctions can be recombined, the refinement can be undone (reversibility), and all operations

preserve admissibility.



\subsection{Interference Closure}

Coherence forces a closure property: if two refinement paths lead to the same physical situation,

their bookkeeping must agree. Consider a distinction d that can be refined two ways: Path 1: d $\to$

{d_1, d_2}, Path 2: d $\to$ {d_1′, d_2′}. If these paths can be coherently recombined, the

bookkeeping structure must be closed under such recombinations. This is interference closure.



\subsection{Linearity}

\textbf{Theorem. Interference closure forces linear structure on the space of coherent realizations.}

The space of coherent bookkeeping states forms a vector space over a division algebra.



\section{Inner Products and Complex Structure}



\subsection{Admissibility-Invariant Overlap}

A similarity measure between coherent realizations must exist that is invariant under admissible

reversible transformations and distributive over refinements. This forces a sesquilinear form: an

inner product.



\subsection{Why Complex Numbers?}

By Solèr’s theorem, the only infinite-dimensional orthomodular spaces over a division algebra are

over R, C, or H. Physical continuity requirements (continuous one-parameter families of reversible

transformations in 2D subspaces) select C. The phase e^{i$\theta$} parameterizes the unique continuous

family of admissibility-preserving rotations in a 2D coherent subspace.



\subsection{Hilbert Space}

Completing the inner product space yields a Hilbert space H. This is the mathematical home of

quantum states: the space of coherent bookkeeping under finite enforceability.



\section{Composition and Tensor Products}



\subsection{Joint Enforceability}

Papers 1--2 established that composition can fail (non-closure). But when composition

succeeds---when two systems are jointly enforceable---what structure does the composite have?



\subsection{Bilinear Composition}

Let A and B be systems with Hilbert spaces H_A and H_B. If they are jointly enforceable, their

composite must have a Hilbert space H_AB. Refinement composability requires that refining A

then composing with B gives the same result as composing then refining. This forces bilinear

composition: $\otimes$: H_A $\times$ H_B $\to$ H_AB.



\subsection{Tensor Product}

\textbf{Theorem. The minimal composite space satisfying bilinear composition and closure under}

admissible recombination is the tensor product H_A $\otimes$ H_B.



\subsection{Entanglement}

Most states in H_A $\otimes$ H_B are not product states. These are entangled states: joint realizations

that cannot be factored into independent realizations. Entanglement is not mysterious in this

framework. It is the generic representation of joint enforceability when the interaction functional is

nonzero. Entangled states encode correlations that require coordination across interfaces.



\subsection{Conditional Composition (Paper 2 Link)}

The tensor product exists if and only if the composite is admissible: A $\otimes$ B exists ⇔ E_$\Gamma$(A $\\cup$ B) $\leq$

C_$\Gamma$. This is non-closure at the level of quantum representations: the familiar assumption that

composition always works is valid only within the admissible region. At saturation (horizons,

measurement), the tensor product description itself breaks down.



\section{The Born Rule}



\subsection{Refinement Invariance}

Probability assignments must be refinement-invariant: coarse-graining and fine-graining must be

\[consistent. If outcome E can be refined into E_1 $\\cup$ E_2, then p(E) = p(E_1) + p(E_2). This is\]

additivity over disjoint outcomes.



\subsection{Admissibility Invariance}

Probability assignments must be admissibility-invariant: they must not depend on arbitrary

bookkeeping choices. For a state $\rho$ and effect E, p(E|$\rho$) must be invariant under basis changes,

equivalent representations, and admissible reversible transformations.



\subsection{Gleason’s Theorem}

\textbf{Theorem (Gleason). In a Hilbert space of dimension $\geq$ 3, the unique probability measure}

\[satisfying additivity and admissibility-invariance is: p(E|$\rho$) = Tr($\rho$E). For pure states $\rho$ =\]



\[|$\\psi$$\\psi$| and projections E = |$\phi$$\phi$|: p = |$\\psi$|$\phi$|². This is the Born rule. [Import: Gleason\]

1957]



\subsection{The Born Rule Is Forced}

The Born rule is not a postulate. It is the unique admissibility-invariant probability assignment in

quantum-admissible regimes. The “mystery” of why probability is |$\\psi$|² dissolves: it is the only

assignment compatible with the structure forced by admissibility.



\section{Dynamics}



\subsection{Reversible Evolution: Unitarity}

\textbf{Theorem. Admissibility-preserving reversible transformations are unitary (or antiunitary).}

\[Preservation of the inner product means Uv, Uw = v, w for all v, w. Physical continuity\]

excludes antiunitary except for time-reversal. Continuous reversible evolution is therefore

\[generated by self-adjoint operators: U(t) = e^{-iHt}. This is the Schrödinger equation\]

structure.



\subsection{General Evolution: CPTP Maps}

Not all evolution is reversible. Record formation, interaction with environments, and capacity

commitment produce irreversible changes. Admissibility-preserving transformations that remain

admissibility-preserving under extension by ancillas are completely positive. With normalization

(trace preservation), admissibility-preserving evolution is described by CPTP maps.



\section{Measurement as Regime Transition}



\subsection{The Standard Problem}

Measurement is the central conceptual difficulty of quantum mechanics. The standard formalism

requires a “collapse postulate” that violates unitarity and lacks a dynamical mechanism. Within the

admissibility framework, measurement requires no postulate. It is a regime transition from

quantum description to record-locked description.



\subsection{Before and After Measurement}

Before measurement: multiple outcomes are jointly admissible, quantum description applies,

superposition is meaningful. After measurement: records correlate with one outcome, capacity is

committed, alternatives contradicting the record are inadmissible. The transition is not mysterious.

It is the same capacity-locking process described in Paper 3 (LEDGERS), now applied to

quantum-coherent distinctions.



\subsection{Record-Locking Channel Formalization (T9)}

The February 2026 release provides an explicit algebraic model for measurement via

\[record-locking channels. The interface algebra decomposes as M = M_sys $\otimes$ Z_R, where M_sys\]

is the system algebra and Z_R is a classical log algebra with projections {p_s} indexed by log

strings.



\[Construction. Each enforcement operation E_i acts as: E_i = (E_i^{sys} $\otimes$ id)  (id $\otimes$ $\alpha$_i),\]

\[where $\alpha$_i is an irreversible append map: $\alpha$_i(p_s) = p_{(s,i)}. The append extends the\]

classical log by one symbol, recording which enforcement was applied.



\textbf{Theorem T9 (L3-$\mu$) [P]. For k enforcement operations applied in all k! orderings, the resulting}

CP maps are pairwise distinct: for $\pi$ $\neq$ $\sigma$, E_$\pi$ $\neq$ E_$\sigma$. The log projections p_$\pi$ and p_$\sigma$ are

orthogonal. Therefore |H_k / $\\equiv$| = k! inequivalent histories.



This resolves the measurement problem in the following sense: measurement is not collapse but

record-locking. The “collapse” is the transition from the coherent regime (where M_sys is the

relevant algebra) to the record-locked regime (where Z_R projections determine which outcome

has been logged). The log is irreversible because record-lock capacity constraints prevent undoing

the append (T0.4′ certificate, Paper 2).



\subsection{Decoherence as Capacity Commitment}

Decoherence is the progressive commitment of capacity to environmental records. As the

environment logs interaction results, the log algebra Z_R grows and the coherent sector M_sys

effectively shrinks. When enough capacity is committed to environmental records, the quantum

description ceases to be applicable---not because quantum mechanics is wrong, but because its

regime conditions (reversible refinement) no longer hold.



\section{Connection to Matter Content}



\subsection{Quantum Structure Constrains Particle Content}

The quantum representation derived in this paper does not float freely above the constraint

structure. It is jointly constrained with the gauge and matter content derived in Papers 1 and 4:



Quantum Structure               Capacity Constraint                 Consequence



Tensor product H_A $\otimes$ H_B Conditional on admissibility               Composition fails at horizons



Unitary symmetry group          Gauge selection (T_gauge)           SU(3)$\times$SU(2)$\times$U(1) uniquely selected



Matter representations          Anomaly cancellation (T5)           Chiral fermion content forced



\[Generation structure            Capacity saturation (T4F)           N_gen = 3 exactly\]



\[Born rule probabilities         Gleason + admissibility             p = |$\\psi$|$\phi$|² uniquely\]



The quantum representation is not free to have arbitrary matter content. The same capacity

architecture that forces Hilbert space also forces three generations of chiral fermions in the

Standard Model gauge group. This is a qualitatively new feature: previous derivations of quantum

structure (Hardy, Chiribella et al., etc.) do not constrain particle content.



\subsection{Particles Require Quantum Structure}

The enforcement potential V($\Phi$) derived in Paper 3 (Section 5.3) shows that particle-like excitations

require the quantum structure derived here. The binding well at $\Phi$/C $\approx$ 0.81 has mass gap d²V/d$\Phi$²

> 0, but no classical soliton localizes: particles are quantum modes of the enforcement field, not

classical cost minima. This closes the conceptual gap between the constraint architecture and

observable particle physics.



\section{Summary}



\subsection{What Is Derived}



Structure                    Source                                     Status



Linear space                 Interference closure                       [P] (from bilinearity regime)



Inner product                Admissibility-invariant overlap            [P] (from regime)



Complex numbers              Phase admissibility + Solèr                [P] + [I]



Hilbert space                Completion of above                        [P]



Tensor products              Compositional closure                      [P] (conditional on admissibility)



Born rule                    Gleason’s theorem                          [P] + [I: Gleason 1957]



Unitarity                    Admissibility-preserving reversible maps [P]



CPTP maps                    Admissibility-preserving general maps      [P]



k! history sectors           Record-locking append maps (T9)            [P]



Measurement                  Regime transition to record-lock           [P_structural]



\subsection{What Is NOT Derived}

Specific Hamiltonians (dynamics depends on physical content). Coupling constants (regime

parameters, but see Paper 4 Section 7 for sin²$\theta$_W). Particle masses beyond the structural hierarchy

(Yukawa ratios are regime parameters).



\subsection{The Status of Quantum Mechanics}

Quantum mechanics is: necessary in bilinear, unsaturated regimes; exact within its domain of

validity; conditional on regime assumptions; limited at saturation, measurement, and horizons. It is

not fundamental in the sense of being the bottom layer. It is a representation---the unique

representation compatible with coherent admissibility bookkeeping in appropriate regimes.



Appendix A: Regime Assumptions

Assumption                Content                                       Where It Fails



Interference closure      Admissible distinctions can be refined/recombined/reversed

Record-locked histories; horizons



Overlap invariance        Similarity measure invariant under admissible

Nearreversible

saturationtransforms

boundaries



Phase admissibility       Continuous 1-parameter reversible families

Discrete

exist inregimes;

2D subspaces

quantized interfaces



Compositional closure Refinement and composition commute forAt

jointly

capacity

enforceable

saturation;

systems

horizon crossing



Bilinearity               Interaction functional I_$\Gamma$ is approximatelyStrong-coupling;

bilinear         near-saturation



Where these regime assumptions fail, quantum mechanics ceases to apply---but the underlying

admissibility constraints (Papers 1--4) continue to hold. The constraint structure is more

fundamental than the quantum representation.



Appendix B: Record-Locking Channel Details

The T9/L3-$\mu$ proof constructs the record-locking model explicitly:



\[(1) Interface algebra: M = M_sys $\otimes$ Z_{$\leq$k}, where Z_{$\leq$k} is the classical log algebra with\]

projections {p_s} for s $\\in$ Σ^m (m $\leq$ k), Σ = {1, ..., k}.



\[(2) Append maps: $\alpha$_i(p_s) = p_{(s,i)} for |s| < k; $\alpha$_i(p_s) = p_s for |s| = k (saturation). Key\]

property: $\alpha$_i(p_s) · $\alpha$_j(p_s) = 0 for i $\neq$ j (ranges disjoint).



\[(3) Enforcement channels: E_i = (E_i^{sys} $\otimes$ id)  (id $\otimes$ $\alpha$_i). Applying E_i refines the\]

system and appends ‘i’ to the irreversible log.



\[(4) For permutation $\pi$: E_{$\pi$(1)}  ...  E_{$\pi$(k)} maps X = 1_sys $\otimes$ p_∅ to 1_sys $\otimes$\]

p_{($\pi$(1),...,$\pi$(k))} = 1_sys $\otimes$ p_$\pi$.



(5) For $\pi$ $\neq$ $\sigma$: p_$\pi$ $\neq$ p_$\sigma$ and p_$\pi$ · p_$\sigma$ = 0. Therefore |H_k/$\\equiv$| = k!.



\textbf{Corollary: dim(Z_R) $\geq$ k! and $\Delta$S_records $\geq$ log(k!) $\sim$ k log k. The entropy of measurement is}

structural, not thermodynamic.



Appendix C: Changelog from January 2026 Release

Change                              Section      Details



Record-locking channels             7.3          T9/L3-$\mu$: explicit algebraic model for measurement



k! history sectors                  7.3          Factorial counting from append maps (proved, not assumed)



Decoherence section                 7.4          Decoherence as progressive capacity commitment to Z_R



Matter content connection           8            New section: quantum structure jointly constrains particle content



Particles require QM                8.2          V($\Phi$) mass gap + no classical solitons $\to$ quantum modes



Derived table expanded              9.1          Added k! histories and measurement with epistemic tags



Change                              Section       Details



Regime table added                  App. A        Explicit table of regime assumptions with failure modes



Forward references                  9             Added Paper 7 (ACTION) and Paper 4 Section 7 (sin²$\theta$_W)



Date updated                        Header        1/30/2026 $\to$ 2/8/2026




\end{document}
